\scanpage{284}
\mspage{26}

La composition des morphismes est la composition des foncteurs entre catégories $I_{\unsure{?|\tau}}$. Ainsi on a un foncteur \uline{pleinement fidèle}$^{*}$ canonique
% slanted
\margin{$*$ faux, cf \unsure{annotation marginale|annotations marginales} page précédente --- on ne prend pas \uline{tous} les foncteurs \unsure{dans|de} $\mathrm{Cat}$ entre $I_\tau$ et $I_{\tau'}$.}
% slanted
\margin{foncteur \uline{réalisation} d'un type (combinatoire) d'intervalle}
\[ \tau\longmapsto I_\tau\colon\quad \underline{\mathrm{Ty}}\xrightarrow{\unsure{\rho|\varphi}}(\mathrm{Cat}), \]
par quoi la catégorie $I_\tau$ \deleted{\unsure{pourrait}} s'identifie à une sous-catégorie pleine de $(\mathrm{Cat})$ (celle des $I_\tau$). Mais je préfère ne pas identifier $\tau\in\mathrm{Ty}$ (donnée combinatoire) avec $I_\tau$ (donnée catégorique).

Soit $X$ une \inserted{(petite ?)} catégorie. On a un contra-foncteur
\[ \left\{\begin{aligned} S&\longmapsto\underline{\mathrm{Hom}}(S,X)\\ (\mathrm{Cat})^{\mathrm{o}}&\mathbin{\text{---}}(\mathrm{Cat})\end{aligned}\right. \]
défini par $X$, d'où en composant avec $\rho$ un contrafoncteur
\diagram{\begin{tikzcd}[ampersand replacement=\&]\underline{\mathrm{Ty}}^{\mathrm{o}} \arrow[r,"{\mathrm{CH}_X}"] \& \mathrm{Cat}\end{tikzcd}}

C'est le contrafoncteur
\[ \left\{\begin{aligned}
\tau&\longmapsto\underline{\mathrm{Hom}}(I_\tau,X)\overset{\mathrm{def}}{=}\underline{\mathrm{Ch}}_\tau(X).\\
&\text{i.e.}\quad \mathrm{CH}_X(\tau)=\underline{\mathrm{Ch}}_\tau(X)
\end{aligned}\right. \]

Mais un contra-foncteur sur une petite catégorie $\underline{\mathrm{Ty}}$ donne lieu à une \uline{catégorie fibrée} (et même scindée) sur celle-ci.
% archive sheet number: 283
